\chapter{2020年浙江卷}

\section{单选题}

\begin{problem}
已知集合 \(P=\{x\mid 1<x<4\}\),\(Q=\{x\mid 2<x<3\}\),则 \(P\cap Q=\)(\quad)

\choices
    {\(\{x\mid 1<x\le 2\}\)}
    {\(\{x\mid 2<x<3\}\)}
    {\(\{x\mid 3\le x<4\}\)}
    {\(\{x\mid 1<x<4\}\)}
\end{problem}

\begin{problem}
已知 \(a\in\R\),若 \(a-1+(a-2)\i\)(\(\i\) 为虚数单位)是实数,则 \(a=\)(\quad)

\choices
    {\(1\)}
    {\(-1\)}
    {\(2\)}
    {\(-2\)}
\end{problem}

\begin{problem}
若实数 \(x\),\(y\) 满足约束条件

\(\begin{cases}
x-3y+1\le0,\\
x+y-3\ge0,
\end{cases}\)

则 \(z=x+2y\) 的取值范围是(\quad)

\choices
    {\((-\infty,4]\)}
    {\([4,+\infty)\)}
    {\([5,+\infty)\)}
    {\((-\infty,+\infty)\)}
\end{problem}

\begin{problem}
函数 \(y=x\cos x+\sin x\) 在区间 \([-\pi,\pi]\) 的图象大致为(\quad)

\choices
    {\(\choicebitmap[width=0.15\paperwidth]{img/2020/zhejiang/q4_choice_a.png}\)}
    {\(\choicebitmap[width=0.15\paperwidth]{img/2020/zhejiang/q4_choice_b.png}\)}
    {\(\choicebitmap[width=0.15\paperwidth]{img/2020/zhejiang/q4_choice_c.png}\)}
    {\(\choicebitmap[width=0.15\paperwidth]{img/2020/zhejiang/q4_choice_d.png}\)}
\end{problem}

\begin{problem}
某几何体的三视图(单位:\(\mathrm{cm}\))如图所示,则该几何体的体积(单位:\(\mathrm{cm}^3\))是(\quad)

\begin{center}
\bitmapfigure[max width=0.46\linewidth]{img/2020/zhejiang/q5_fig1.png}
\end{center}

\choices
    {\(\dfrac{7}{3}\)}
    {\(\dfrac{14}{3}\)}
    {\(3\)}
    {\(6\)}
\end{problem}

\begin{problem}
已知空间中不过同一点的三条直线 \(m\),\(n\),\(l\),则``\(m\),\(n\),\(l\) 在同一平面''是``\(m\),\(n\),\(l\) 两两相交''的(\quad)

\choices
    {充分不必要条件}
    {必要不充分条件}
    {充分必要条件}
    {既不充分也不必要条件}
\end{problem}

\begin{problem}
已知等差数列 \(\{a_n\}\) 的前 \(n\) 项和为 \(S_n\),公差 \(d\neq0\),且 \(\dfrac{a_1}{d}\le1\). 记 \(b_1=S_2\),\(b_{n+1}=S_{2n+2}-S_{2n}\)(\(n\in\N^*\)),下列等式不可能成立的是(\quad)

\choices
    {\(2a_4=a_2+a_6\)}
    {\(2b_4=b_2+b_6\)}
    {\(a_4^2=a_2a_8\)}
    {\(b_4^2=b_2b_8\)}
\end{problem}

\begin{problem}
已知点 \(O(0,0)\),\(A(-2,0)\),\(B(2,0)\). 设点 \(P\) 满足 \(|PA|-|PB|=2\),且 \(P\) 为函数 \(y=3\sqrt{4-x^2}\) 图象上的点,则 \(|OP|=\)(\quad)

\choices
    {\(\dfrac{\sqrt{22}}{2}\)}
    {\(\dfrac{4\sqrt{10}}{5}\)}
    {\(\sqrt7\)}
    {\(\sqrt{10}\)}
\end{problem}

\begin{problem}
已知 \(a\),\(b\in\R\) 且 \(ab\neq0\),若 \((x-a)(x-b)(x-2a-b)\ge0\) 在 \(x\ge0\) 上恒成立,则(\quad)

\choices
    {\(a<0\)}
    {\(a>0\)}
    {\(b<0\)}
    {\(b>0\)}
\end{problem}

\begin{problem}
设集合 \(S\),\(T\),\(S\subseteq\N^*\),\(T\subseteq\N^*\),\(S\),\(T\) 中至少有两个元素,且 \(S\),\(T\) 满足:

\begin{circlelist}
\item 对于任意 \(x\),\(y\in S\),若 \(x\ne y\),都有 \(xy\in T\);

\item 对于任意 \(x\),\(y\in T\),若 \(x<y\),则 \(\dfrac{y}{x}\in S\).
\end{circlelist}

下列命题正确的是(\quad)

\choices
    {若 \(S\) 有 4 个元素,则 \(S\cup T\) 有 7 个元素}
    {若 \(S\) 有 4 个元素,则 \(S\cup T\) 有 6 个元素}
    {若 \(S\) 有 3 个元素,则 \(S\cup T\) 有 4 个元素}
    {若 \(S\) 有 3 个元素,则 \(S\cup T\) 有 5 个元素}
\end{problem}

\section{填空题}

\begin{problem}
已知数列 \(\{a_n\}\) 满足 \(a_n=\dfrac{n(n+1)}{2}\),则 \(S_3=\fillinblank{}\).
\end{problem}

\begin{problem}
设 \((1+2x)^5=a_1+a_2x+a_3x^2+a_4x^3+a_5x^4+a_6x^5\),则 \(a_5=\fillinblank{}\);\(a_1+a_2+a_3=\fillinblank{}\).
\end{problem}

\begin{problem}
已知 \(\tan\theta=2\),则 \(\cos2\theta=\fillinblank{}\);\(\tan\left(\theta-\dfrac{\pi}{4}\right)=\fillinblank{}\).
\end{problem}

\begin{problem}
已知圆锥展开图的侧面积为 \(2\pi\),且为半圆,则底面半径为\fillinblank{}.
\end{problem}

\begin{problem}
设直线 \(l:y=kx+b\)(\(k>0\)),圆 \(C_1:x^2+y^2=1\),\(C_2:(x-4)^2+y^2=1\),若直线 \(l\) 与 \(C_1\),\(C_2\) 都相切,则 \(k=\fillinblank{}\);\(b=\fillinblank{}\).
\end{problem}

\begin{problem}
一个盒子里有 1 个红,1 个绿,2 个黄四个相同的球,每次拿一个,不放回,拿出红球即停,设拿出黄球的个数为 \(\xi\),则 \(P(\xi=0)=\fillinblank{}\);\(E(\xi)=\fillinblank{}\).
\end{problem}

\begin{problem}
设 \(\bs e_1\),\(\bs e_2\) 为单位向量,满足 \(|2\bs e_1-\bs e_2|\le\sqrt2\),\(\bs a=\bs e_1+\bs e_2\),\(\bs b=3\bs e_1+\bs e_2\),设 \(\bs a\),\(\bs b\) 的夹角为 \(\theta\),则 \(\cos^2\theta\) 的最小值为\fillinblank{}.
\end{problem}

\section{解答题}

\begin{problem}
在锐角 \(\triangle ABC\) 中,角 \(A\),\(B\),\(C\) 的对边分别为 \(a\),\(b\),\(c\),且 \(2b\sin A=\sqrt3\,a\).

\begin{enumerate}
\item 求角 \(B\);

\item 求 \(\cos A+\cos B+\cos C\) 的取值范围.
\end{enumerate}
\end{problem}

\begin{problem}
如图,三棱台 \(DEF-ABC\) 中,面 \(ADFC\perp\) 面 \(ABC\),\(\angle ACB=\angle ACD=45^\circ\),\(DC=2BC\).

\begin{enumerate}
\item 证明:\(EF\perp DB\);

\item 求 \(DF\) 与面 \(DBC\) 所成角的正弦值.
\end{enumerate}

\begin{center}
\bitmapfigure[max width=0.36\linewidth]{img/2020/zhejiang/q19_fig1.png}
\end{center}
\end{problem}

\begin{problem}
已知数列 \(\{a_n\}\),\(\{b_n\}\),\(\{c_n\}\) 中,\(a_1=b_1=c_1=1\),\(c_{n+1}=a_{n+1}-a_n\),\(c_{n+1}=\dfrac{b_n}{b_{n+2}}\,c_n\)(\(n\in\N^*\)).

\begin{enumerate}
\item 若数列 \(\{b_n\}\) 为等比数列,且公比 \(q>0\),且 \(b_1+b_2=6b_3\),求 \(q\) 与 \(a_n\) 的通项公式;

\item 若数列 \(\{b_n\}\) 为等差数列,且公差 \(d>0\),证明:\(c_1+c_2+\cdots+c_n<1+\dfrac1d\),\(n\in\N^*\).
\end{enumerate}
\end{problem}

\begin{problem}
如图,已知椭圆 \(C_1:\dfrac{x^2}{2}+y^2=1\),抛物线 \(C_2:y^2=2px\)(\(p>0\)),点 \(A\) 是椭圆 \(C_1\) 与抛物线 \(C_2\) 的交点,过点 \(A\) 的直线 \(l\) 交椭圆 \(C_1\) 于点 \(B\),交抛物线 \(C_2\) 于 \(M\)(\(B\),\(M\) 不同于 \(A\)).

\begin{enumerate}
\item 若 \(p=\dfrac{1}{16}\),求抛物线 \(C_2\) 的焦点坐标;

\item 若存在不过原点的直线 \(l\) 使 \(M\) 为线段 \(AB\) 的中点,求 \(p\) 的最大值.
\end{enumerate}

\begin{center}
\bitmapfigure[max width=0.38\linewidth]{img/2020/zhejiang/q21_fig1.png}
\end{center}
\end{problem}

\begin{problem}
已知 \(1<a\le2\),函数 \(f(x)=\e^x-x-a\),其中 \(\e=2.71828\cdots\) 为自然对数的底数.

\begin{enumerate}
\item 证明:函数 \(y=f(x)\) 在 \((0,+\infty)\) 上有唯一零点;

\item 记 \(x_0\) 为函数 \(y=f(x)\) 在 \((0,+\infty)\) 上的零点,证明:

\begin{enumerate}
\item \(\sqrt{a-1}\le x_0\le\sqrt{2(a-1)}\);
\item \(x_0f(\e^{x_0})\ge(\e-1)(a-1)a\).
\end{enumerate}
\end{enumerate}
\end{problem}

